\chapter{Pipeline DataOps}

\section{Ingestion avec dlt}

Le pipeline d'ingestion est implémenté dans \repoPath{fraud-detection/dataops/ingestion.py}. Il charge le fichier Kaggle, applique le contrat via \repoPath{assert_creditcard_contract}, puis exécute un pipeline dlt vers DuckDB.

La table cible est \repoPath{raw.creditcard_transactions}. Le mode \repoPath{replace} est utilisé pour garantir une reconstruction reproductible de l'état raw à partir de la source.

\section{Stockage DuckDB}

DuckDB est utilisé comme entrepôt local, léger et reproductible. Il répond au besoin pédagogique de disposer d'un stockage analytique sans dépendance à une infrastructure cloud. Le chemin configuré est \repoPath{data/warehouse/fraud.duckdb}.

\section{Transformations dbt}

Le projet dbt est placé dans \repoPath{fraud-detection/dbt_fraud}. Il contient deux niveaux principaux:

\begin{itemize}
    \item \repoPath{stg_transactions}: table nettoyée, typée et dédupliquée;
    \item \repoPath{fraud_features}: table orientée ML contenant les variables originales et les features dérivées.
\end{itemize}

\subsection{Staging}

Le modèle \repoPath{stg_transactions.sql} sélectionne les lignes valides, caste les types et supprime les doublons. Les règles principales sont:

\begin{itemize}
    \item \repoPath{Time}, \repoPath{Amount} et \repoPath{Class} non nuls;
    \item \repoPath{Class} limitée aux valeurs 0 et 1;
    \item \repoPath{Amount} positif ou nul;
    \item variables PCA castées en double.
\end{itemize}

\subsection{Mart de Features}

Le modèle \repoPath{fraud_features.sql} produit des variables utiles au modèle:

\begin{itemize}
    \item heure transactionnelle;
    \item période de la journée;
    \item transformation logarithmique du montant;
    \item binning du montant;
    \item moyenne et écart-type des variables PCA;
    \item score de risque composite;
    \item interactions \repoPath{amount_x_v14} et \repoPath{amount_x_v17}.
\end{itemize}

\section{Orchestration Dagster}

Dagster matérialise les assets suivants:

\begin{enumerate}
    \item \repoPath{download_dataset};
    \item \repoPath{ingest_data};
    \item \repoPath{dbt_transform};
    \item \repoPath{train_model};
    \item \repoPath{evaluate_model};
    \item \repoPath{monitor_model}.
\end{enumerate}

Cette orchestration relie le DataOps au MLOps: la donnée validée alimente dbt, puis le pipeline d'entraînement et les métriques de monitoring.

\section{Reproductibilité}

La reproductibilité est assurée par:

\begin{itemize}
    \item une configuration centralisée dans \repoPath{config/settings.py};
    \item un \repoPath{RANDOM_STATE} utilisé dans les splits, SMOTE et modèles;
    \item des artefacts sauvegardés sous \repoPath{models/};
    \item des commandes documentées dans le README;
    \item Docker Compose pour aligner les services d'exécution.
\end{itemize}
